% Resumo (NBR 6028: um parágrafo, 150-500 palavras, sem fórmulas).
% Todo número rastreia aos capítulos 4-7 (ver comentários de proveniência
% nas fontes dos capítulos).
A escolha de modelos para dados tabulares foi guiada, por uma década, pela
máxima de que o \emph{gradient boosting} sempre vence; a chegada dos
\emph{foundation models} tabulares, redes pré-treinadas que predizem
sem treinamento na tarefa, tornou essa orientação obsoleta sem lhe dar
um substituto apoiado em evidência neutra. Esta dissertação ataca o
problema em três frentes encadeadas. Primeiro, um benchmark multicritério
avalia 14 modelos de três famílias (\emph{gradient boosting}, aprendizado
profundo e \emph{foundation models}) em 18 conjuntos de dados públicos do
OpenML (250 das 252 células possíveis, com duas exclusões estruturais
declaradas), sob \emph{hold-out} intocado pela seleção e
inferência estatística frequentista e bayesiana. No núcleo de 11 modelos
da geração anterior a 2026, o topo da acurácia é um empate estatístico (8 dos 10 modelos não
líderes são praticamente equivalentes ao líder sob região de equivalência
prática de um ponto percentual de AUC, no critério de massa posterior
esperada acima de meio, declarado no texto; três de dez no corte
estrito), e a geração 2026 o quebra: o
TabFM, avaliado \emph{zero-shot}, lidera o rank médio nas três tarefas
(2,7 na classificação binária, 1,0 na multiclasse e 1,2 na regressão) e é
o melhor modelo em 13 dos 18 conjuntos, ao custo de quatro a cinco ordens
de magnitude em latência de inferência; as redes de Kolmogorov--Arnold,
testadas de forma independente, terminam no último terço nas três
tarefas. Segundo, um arcabouço de decisão é avaliado por
\emph{leave-one-dataset-out} (validação negativa do roteador
aprendido, avaliação comparativa das políticas fixas): o roteamento
por meta-características não
generaliza (taxa de acerto de 0,11 contra 0,44 do \emph{baseline}
trivial), a política fixa ``\emph{foundation model} primeiro'' é ótima em
\emph{regret} na métrica primária (\emph{regret} mediano de 0,000 na
geração 2026, preservado na política executável de modelo único) e um
experimento registrado de \emph{regret}
multicritério mostra que basta deslocar 10 a 20 por cento do peso da
acurácia para a latência (cruzamento entre os pesos 0,9 e 0,8, na
normalização por postos) para a
recomendação deixar de ser o \emph{foundation model},
justificando quantitativamente o fluxograma de quatro restrições verificáveis a priori que
consolida o arcabouço. Terceiro, a destilação distribucional de
\emph{foundation models}, recorte verificado como inédito,
transfere as curvas de quantis do professor TabPFN para alunos que servem
em microssegundos: a destilação pontual falha de forma robusta em escala
realista de dados, mas a distribucional retém mediana de 19 por cento da
vantagem distribucional do professor com hiperparâmetros herdados
(12 por cento com os dois braços re-ajustados, a estimativa
conservadora), positiva em 16 de 20 conjuntos
(teste de sinal com p-valor de 0,006 na leitura direta e de
aproximadamente 0,019 sob a correção sequencial do desenho em duas
etapas, declarada no texto; 15 de 19, com p-valor de 0,010,
na sensibilidade que exclui o conjunto de alvo binarizado da errata),
vence \emph{baselines}
distribucionais clássicos em 17 de 20 conjuntos no experimento
pré-registrado E1 e se condensa numa regra de decisão prática cujas
pré-condições se verificam em minutos. O conjunto (evidência sob
protocolo único, arcabouço avaliado sob LODO, destilação seletiva)
substitui a
máxima por um procedimento; o que o \emph{foundation model} sabe e os
dados brutos não ensinam ao aluno é a distribuição preditiva, e é isso
que se destila.

\vspace{1cm}
\noindent\textbf{Palavras-chave:} dados tabulares; foundation models;
benchmark; destilação de conhecimento; regressão distribucional.
